% =====================================================================
\section{Fundamentos da Troca de Chaves Pós-Quântica}\label{sec:fund-pqc}
% =====================================================================
Esta seção explica, desde o início, por que um nó de IoT precisa combinar uma
chave, como isso é feito hoje, por que os computadores quânticos ameaçam esse
processo e como funciona o novo padrão do NIST, o ML-KEM, cujas mensagens são
muito maiores.

\subsection{Por que Dois Dispositivos Precisam de uma Chave em Comum}\label{sec:fund-sym}
A maior parte dos dados que um sensor envia é protegida com criptografia
\emph{simétrica}: emissor e receptor possuem a mesma chave secreta, e a mesma
chave serve para cifrar e para decifrar. Cifras simétricas como o AES, ou a
família leve Ascon, padronizada pelo NIST em 2025 para dispositivos
restritos~\cite{sp800232}, são rápidas e baratas o bastante para qualquer
microcontrolador.

A criptografia simétrica deixa uma pergunta em aberto: como os dois
dispositivos passam a ter a mesma chave? Um nó tirado da caixa e instalado em
uma enfermaria nunca ``conheceu'' o ponto de acesso com o qual vai conversar, e
tudo o que um diz ao outro trafega por um enlace que qualquer pessoa na sala
pode observar (Seção~\ref{sec:fund-light}). Enviar a chave às claras a
revelaria. Esse é o \emph{problema da distribuição de chaves}, e ele é
resolvido pela \emph{criptografia de chave pública}. Uma analogia conhecida é a
do cadeado aberto. Suponha que o ponto de acesso queira receber um segredo: ele
envia ao nó um cadeado \emph{aberto} e guarda a única chave. O nó coloca um
segredo recém-escolhido em uma caixa, fecha o cadeado e devolve a caixa.
Qualquer um pode ver o cadeado e a caixa em trânsito, mas só o ponto de acesso
consegue abri-la. Os algoritmos de chave pública usados neste artigo
funcionam, em essência, como esse cadeado.

\subsection{Troca de Chaves Clássica: Diffie--Hellman e X25519}\label{sec:fund-dh}
A troca de chaves usada hoje em quase toda parte é a de Diffie--Hellman. Uma
analogia popular usa tintas. As duas partes começam com uma cor comum e
pública. Cada uma mistura a ela uma cor secreta e envia a mistura à outra. Em
seguida, cada uma acrescenta a sua própria cor secreta ao que recebeu. As duas
terminam com a mesma cor final, a cor comum mais os dois segredos, enquanto um
observador que viu apenas as duas misturas não consegue separar os segredos,
porque ``desmisturar'' tinta é difícil.

Na versão matemática, ``misturar'' é uma operação sobre os pontos de uma curva
elíptica, e ``desmisturar'' é o \emph{problema do logaritmo discreto em curvas
elípticas}, para o qual não se conhece algoritmo clássico eficiente. A variante
mais usada, a X25519~\cite{bernstein2006curve25519,rfc7748}, é compacta: cada
parte envia um único valor de 32 bytes, de modo que a troca completa movimenta
$\val{X25519Bytes}$ bytes. Ela também é barata o bastante para rodar em um
microcontrolador. Há duas décadas esse é o modo padrão de pequenos dispositivos
combinarem uma chave.

\subsection{A Ameaça Quântica}\label{sec:fund-quantum}
Um \emph{computador quântico} processa informação usando efeitos da mecânica
quântica e, para alguns problemas específicos, pode ser dramaticamente mais
rápido do que qualquer computador comum. Um desses problemas é exatamente o que
protege a troca de Diffie--Hellman. Em 1994, Peter Shor mostrou que um
computador quântico suficientemente grande resolveria com eficiência tanto a
fatoração de inteiros quanto o logaritmo discreto~\cite{shor1997}, o que
quebraria o RSA, o Diffie--Hellman e a X25519.

A criptografia simétrica resiste muito melhor. O melhor ataque quântico
conhecido contra ela, o algoritmo de Grover, apenas acelera a busca exaustiva
da chave por uma raiz quadrada, e o remédio usual, usar chaves mais longas,
como o AES de 256 bits, é simples. O ponto fraco de uma implantação de IoT na
era quântica não é, portanto, a cifração dos dados, e sim a \emph{troca de
chaves} que a prepara.

Não existe hoje um computador quântico grande o bastante para executar o
algoritmo de Shor contra chaves reais, e as estimativas de quando existirá
variam muito. Por que se preocupar agora? Porque o atacante não precisa do
computador quântico no momento do ataque. Ele pode gravar hoje uma sessão
cifrada, armazená-la e decifrá-la quando a máquina estiver disponível: é a
ameaça conhecida como \emph{colher agora, decifrar depois} (\emph{harvest now,
decrypt later}). Mosca~\cite{mosca2018quantum} resume o prazo resultante em uma
desigualdade simples. Seja $x$ o número de anos durante os quais os dados
precisam permanecer confidenciais, $y$ o número de anos necessários para trocar
a criptografia dos dispositivos já instalados e $z$ o número de anos até
existir um computador quântico de grande porte. Então, se
\begin{equation}
x + y > z,
\label{eq:mosca}
\end{equation}
os dados coletados hoje já estão expostos. Em IoT, os dois termos do lado
esquerdo da Eq.~\eqref{eq:mosca} são grandes: dados médicos e industriais
muitas vezes precisam permanecer privados por muitos anos, e dispositivos
embarcados, uma vez instalados, podem ficar uma década sem atualização. É por
isso que a migração para a criptografia pós-quântica é urgente justamente para
o tipo de dispositivo considerado neste artigo.

\emph{Criptografia pós-quântica} significa algoritmos comuns, que rodam em
processadores comuns, projetados para resistir a ataques de computadores
quânticos, além dos clássicos. Ela não exige hardware quântico, apenas software
novo. O NIST conduziu uma competição aberta, de vários anos, para selecionar
esses algoritmos; a Tabela~\ref{tab:nist} resume seus marcos.

\begin{table}[!htb]
\caption{Principais marcos da padronização pós-quântica do NIST.}
\label{tab:nist}
\centering
\small
\begin{tabular}{ll}
\toprule
ano & marco\\
\midrule
2016 & o NIST abre a chamada pública por algoritmos pós-quânticos\\
2022 & CRYSTALS-Kyber é selecionado para estabelecimento de chaves\\
2024 & publicação do FIPS 203 (ML-KEM, derivado do Kyber)~\cite{fips203}\\
     & e do FIPS 204 (ML-DSA), para assinaturas~\cite{fips204}\\
2025 & HQC é selecionado como segundo KEM, baseado em códigos~\cite{nisthqc2025};\\
     & SP 800-227, com recomendações de uso de KEMs~\cite{sp800227}\\
\bottomrule
\end{tabular}
\end{table}

\subsection{Mecanismos de Encapsulamento de Chaves}\label{sec:fund-kem}
O substituto pós-quântico do Diffie--Hellman é apresentado na forma de um
\emph{mecanismo de encapsulamento de chaves} (KEM, \emph{key-encapsulation
mechanism})~\cite{cramer2003kem,sp800227}. Um KEM é o cadeado da
Seção~\ref{sec:fund-sym} em versão precisa. Ele é formado por três algoritmos:
\begin{itemize}
\item $\textsc{KeyGen}() \to (\ek, \dk)$ fabrica um cadeado novo e a sua chave:
      uma \emph{chave de encapsulamento} pública $\ek$ e uma \emph{chave de
      desencapsulamento} privada $\dk$;
\item $\textsc{Encaps}(\ek) \to (K, \ct)$ escolhe um segredo aleatório novo $K$
      e o tranca: o \emph{texto cifrado} $\ct$ é a caixa fechada;
\item $\textsc{Decaps}(\dk, \ct) \to K$ abre a caixa e recupera $K$.
\end{itemize}
Em uma troca de chaves, uma parte executa \textsc{KeyGen} e envia $\ek$; a
outra executa \textsc{Encaps}, guarda $K$ e devolve $\ct$; a primeira executa
\textsc{Decaps} e passa a ter o mesmo $K$. Duas diferenças em relação ao
Diffie--Hellman serão importantes adiante. As duas mensagens são objetos
diferentes, de tamanhos diferentes, e as duas partes realizam quantidades
diferentes de trabalho. Ambas afetam o comportamento da troca em um enlace de
capacidade limitada.

Um KEM precisa continuar seguro mesmo contra um atacante capaz de injetar
textos cifrados de sua autoria e observar como o receptor reage. O nome formal
dessa propriedade é segurança IND-CCA2. Os KEMs baseados em reticulados a
obtêm com uma técnica chamada transformada de
Fujisaki--Okamoto~\cite{fujisaki1999,hofheinz2017fo}: ao abrir uma caixa, o
\textsc{Decaps} refaz a caixa a partir do que encontrou dentro dela e verifica
se o resultado é idêntico à caixa recebida. Se um único bit for diferente, a
caixa foi adulterada ou danificada.

O que o \textsc{Decaps} faz em seguida é uma decisão de projeto sutil, mas
importante. A escolha óbvia, devolver um erro, seria um equívoco: um atacante
que possa enviar muitos textos cifrados modificados e observar quais produzem
erro aprende algo com cada resposta, e ataques exatamente desse tipo já
quebraram sistemas reais~\cite{bleichenbacher1998}. O ML-KEM usa, em vez disso,
a \emph{rejeição implícita}: quando a verificação falha, ele devolve
silenciosamente um valor de aparência aleatória, calculado a partir de uma
semente secreta e do texto cifrado~\cite{fips203}. O atacante não aprende nada,
e as duas partes legítimas simplesmente ficam com chaves diferentes, o que
percebem na primeira tentativa de usá-las. Em um enlace ruidoso, essa é também
a última linha de defesa contra erros de transmissão que escapem ao código de
verificação da Seção~\ref{sec:fund-arq}.

Um KEM, sozinho, não confirma a identidade de ninguém; ele apenas estabelece um
segredo. Por isso, o NIST recomenda alimentar uma \emph{função de derivação de
chaves} com esse segredo e com as mensagens trocadas~\cite{sp800227}. Uma
função de derivação de chaves é construída a partir de uma \emph{função de
hash}, um algoritmo que transforma qualquer entrada em uma ``impressão
digital'' de tamanho fixo, de modo que entradas diferentes produzem, na
prática, impressões diferentes. Derivar a chave de sessão do segredo do KEM e
de uma impressão digital de toda a conversa amarra a chave àquela troca
específica, de forma que uma troca adulterada gera uma chave inútil.

\subsection{O ML-KEM e a Dificuldade das Equações com Ruído}\label{sec:fund-mlkem}
O ML-KEM (\emph{Module-Lattice-based KEM}), padronizado como FIPS 203 e
derivado do CRYSTALS-Kyber~\cite{fips203,bos2018kyber}, apoia-se em um problema
fácil de enunciar e, até onde se sabe, difícil de resolver mesmo com um
computador quântico: resolver um sistema de equações lineares em que cada
equação recebeu um pequeno erro.

Um exemplo de brinquedo transmite a ideia. Trabalhando módulo $17$, suponha
que lhe informem
\begin{equation}
3s_1 + 5s_2 \equiv 2, \qquad 7s_1 + 2s_2 \equiv 9 \pmod{17}.
\label{eq:toy}
\end{equation}
O sistema da Eq.~\eqref{eq:toy} pode ser resolvido por eliminação, exatamente
como na álgebra escolar, e dá $s_1=2$, $s_2=6$. Suponha agora que cada lado
direito tenha sido alterado por um pequeno valor aleatório, $-1$, $0$ ou $+1$,
e que você não saiba qual. A eliminação multiplica e subtrai as equações, e
multiplica e subtrai junto os erros ocultos; depois de alguns passos, os erros
cresceram até o tamanho dos próprios números, e a resposta perde o sentido. Com
duas incógnitas ainda é possível testar todas as possibilidades, mas no
ML-KEM-512 o segredo tem $512$ incógnitas, muito além de qualquer busca
exaustiva. Esse é o problema do \emph{aprendizado com erros} (LWE,
\emph{learning with errors}), introduzido por Regev~\cite{regev2009lwe}; o
ML-KEM usa uma variante estruturada, o LWE em \emph{módulos} (MLWE), que
preserva a dificuldade e torna a aritmética muito mais
rápida~\cite{langlois2015mlwe}.

No ML-KEM, os números são substituídos por polinômios de 256 coeficientes, com
toda a aritmética feita módulo $q=3329$. A chave pública contém uma matriz
aleatória $\mathbf{A}$, de dimensão $k\times k$, formada por esses polinômios,
e um vetor~\cite{fips203}
\begin{equation}
\mathbf{t} = \mathbf{A}\mathbf{s} + \mathbf{e},
\label{eq:mlwe}
\end{equation}
em que o segredo $\mathbf{s}$ e o erro $\mathbf{e}$ têm coeficientes pequenos.
A Eq.~\eqref{eq:mlwe} é o sistema com ruído da Eq.~\eqref{eq:toy} em escala
muito maior: recuperar $\mathbf{s}$ a partir de $(\mathbf{A},\mathbf{t})$ é o
problema difícil. A matriz $\mathbf{A}$ nunca é enviada de fato: os dois lados
a regeneram a partir de uma semente $\rho$ de 32 bytes. A dimensão $k$ define o
nível de segurança. A Tabela~\ref{tab:mlkem} lista os três conjuntos de
parâmetros do padrão, que visam às categorias de segurança 1, 3 e 5 do NIST,
aproximadamente tão difíceis de quebrar quanto o AES-128, o AES-192 e o
AES-256.

\begin{table}[!htb]
\caption{Conjuntos de parâmetros do ML-KEM (FIPS 203). Tamanhos em bytes.
$\delta$ é a probabilidade de uma troca honesta falhar por razões
criptográficas.}
\label{tab:mlkem}
\centering
\small
\begin{tabular}{lcccccc}
\toprule
conjunto & $k$ & $(\eta_1;\eta_2)$ & $(d_u;d_v)$ & $|\ek|$ & $|\ct|$ & $\delta$\\
\midrule
ML-KEM-512  & 2 & (3; 2) & (10; 4) & \phantom{1}800  & \phantom{1}768  & $2^{-138{,}8}$\\
ML-KEM-768  & 3 & (2; 2) & (10; 4) & 1184 & 1088 & $2^{-164{,}8}$\\
ML-KEM-1024 & 4 & (2; 2) & (11; 5) & 1568 & 1568 & $2^{-174{,}8}$\\
\bottomrule
\end{tabular}
\end{table}

Os tamanhos das mensagens decorrem dessa estrutura, e vale a pena deduzi-los,
porque eles determinam todos os resultados deste artigo. A chave de
encapsulamento carrega $\mathbf{t}$, que são $k$ polinômios de 256
coeficientes; cada coeficiente é um número menor que $3329$ e precisa de 12
bits. Somando a semente de 32 bytes, obtém-se~\cite{fips203}
\begin{equation}
|\ek| = \frac{256\times12}{8}\,k + 32 = 384k + 32 \ \text{bytes},
\label{eq:eksize}
\end{equation}
ou seja, $800$ bytes no ML-KEM-512, em que $k=2$. O texto cifrado é formado por
$k+1$ polinômios \emph{comprimidos}, dos quais se guardam apenas $d_u$ ou $d_v$
bits por coeficiente~\cite{fips203}:
\begin{equation}
|\ct| = 32\,(d_u k + d_v) \ \text{bytes},
\label{eq:ctsize}
\end{equation}
o que dá $768$ bytes no ML-KEM-512. As Eqs.~\eqref{eq:eksize}
e~\eqref{eq:ctsize} reproduzem todas as linhas da Tabela~\ref{tab:mlkem}. O
segredo compartilhado $K$ tem sempre $32$ bytes. As probabilidades de falha
$\delta$ da tabela são tão pequenas, abaixo de uma em $10^{41}$, que na prática
uma troca honesta nunca falha por razões criptográficas. Em um enlace ruidoso,
toda falha é uma falha de transmissão.

Embora as chaves sejam grandes, o cálculo é rápido. A multiplicação de
polinômios é feita com a transformada numérica (NTT), uma versão da FFT que
trabalha com números inteiros, e toda a aleatoriedade e todo o hash usam o
SHAKE, da família SHA-3~\cite{fips202}. Em um ARM Cortex-M4, uma geração de
chaves do ML-KEM-512 leva menos de $400\,000$ ciclos, o que, pela
Eq.~\eqref{eq:ciclos}, corresponde a cerca de $2{,}3$\,ms a
$168$\,MHz~\cite{pqm4}.

\subsection{Por que as Trocas de Chaves Pós-Quânticas São Grandes}\label{sec:fund-size}
Colocar os números lado a lado torna visível a tensão central deste artigo.
Uma troca X25519 movimenta $\val{X25519Bytes}$ bytes; uma troca ML-KEM-512
movimenta $\val{MLKEM512Bytes}$, cerca de $25$ vezes mais. A diferença não é uma
ineficiência do projeto: um problema de equações com ruído só é difícil se
contiver muitas equações, e cada equação contribui com coeficientes que
precisam ser transmitidos. A criptografia baseada em reticulados troca tamanho
por velocidade. Ela é maior do que a de curvas elípticas no enlace, mas, como
mostra a Seção~\ref{sec:results}, não é mais lenta para calcular em um
microcontrolador.

Outros KEMs pós-quânticos fazem trocas diferentes. Os esquemas \emph{baseados
em códigos} apoiam-se na dificuldade de decodificar um código corretor de erros
aleatório, problema estudado há mais de quarenta anos. O NIST selecionou o HQC
em 2025 como segundo KEM padrão justamente para que um eventual avanço contra
os reticulados não deixasse o mundo sem um padrão~\cite{nisthqc2025}. O BIKE
foi uma alternativa próxima, e o Classic McEliece é o mais conservador de
todos. Seus tamanhos, listados na Tabela~\ref{tab:kems}, têm formatos
incomuns: o HQC-128 e o BIKE-L1 enviam alguns quilobytes em cada sentido,
enquanto o Classic McEliece envia um texto cifrado minúsculo, de 96 bytes, mas
exige uma chave pública de 255 quilobytes. Em um enlace no qual cada byte custa
tempo de transmissão, esses formatos levam a resultados muito diferentes, que
são o assunto do restante do artigo.
